\subsection{Parameterdarstellung einer Ebene}
\label{subsec:analytic_geometry_parametric_plane_equations}

Ein Aufpunkt und zwei Bewegungsrichtungen erzeugen alle Punkte einer Ebene.
Die lineare Unabhängigkeit der Richtungen stellt sicher, dass dabei wirklich
eine Fläche und nicht nur eine Gerade entsteht.

\begin{DefinitionBox}{Parameterdarstellung einer Ebene}{analytic_geometry_parametric_plane}
  Eine Ebene $E$ durch den Punkt $P$ mit Ortsvektor $\vec p$ und mit zwei
  linear unabhängigen Richtungsvektoren $\vec a,\vec b\in\R^3$ besitzt die
  Parameterdarstellung
  \[
    E:\quad
    \vec x=\vec p+\lambda\vec a+\mu\vec b,
    \qquad \lambda,\mu\in\R.
  \]
  Die Parameter $\lambda$ und $\mu$ steuern unabhängig die Bewegung in den
  beiden Richtungen der Ebene.
\end{DefinitionBox}

\begin{center}
  \CachedTikzFigure{content/resources/chapters/analytic_geometry/tikz/plane_parameter_representation.tex}
\end{center}

\begin{KeyPointBox}{Wann spannen zwei Richtungen eine Ebene auf?}{analytic_geometry_plane_direction_vectors}
  Für eine Ebene müssen $\vec a$ und $\vec b$ linear unabhängig sein:
  \[
    \vec a\ne\vec 0,\qquad
    \vec b\ne\vec 0,\qquad
    \vec b\ne t\vec a\quad\text{für alle }t\in\R.
  \]
  Andernfalls erzeugen beide Parameter nur Punkte auf einer Geraden.
\end{KeyPointBox}

\begin{RemarkBox}[title=Nicht eindeutige Darstellung]
  Wie bei Geraden kann der Aufpunkt durch jeden anderen Punkt der Ebene
  ersetzt werden. Auch jedes andere Paar linear unabhängiger Richtungsvektoren
  in der Ebene beschreibt dieselbe Ebene.
\end{RemarkBox}
